\section{Power Spectrum in Spherical Fourier-Bessel Space: Theory and Application}

This section provides a comprehensive exposition of power spectrum theory in the Spherical Fourier-Bessel (SFB)
decomposition framework, emphasizing the physical meaning of each quantity and the mathematical structure that
emerges when fields are observed on the curved sky over extended redshift ranges.

\subsection{Field Decomposition and the Spherical Fourier-Bessel Basis}

\subsubsection{The homogeneous and isotropic field in SFB coordinates}

Consider a three-dimensional scalar field $f(\mathbf{r})$ defined in Euclidean space with spherical coordinates
$\mathbf{r}=(r,\theta,\phi)$, where $r$ is the radial distance, $\theta$ is the polar angle, and $\phi$ is the azimuthal angle.
In cosmological applications, this field typically represents the matter overdensity $\delta(\mathbf{r})$ or galaxy count
fluctuations.

The field can be expanded in the complete orthonormal basis of spherical Fourier-Bessel functions:
\begin{equation}
f(\mathbf{r}) = \sqrt{\frac{2}{\pi}}
\int_0^\infty dk\,
\sum_{\ell=0}^\infty\sum_{m=-\ell}^{\ell}
f_{\ell m}(k)\,k j_\ell(kr)\,Y_{\ell m}(\theta,\phi),
\label{eq:sfb_expansion}
\end{equation}
where $j_\ell(x)$ are the spherical Bessel functions of the first kind, and $Y_{\ell m}(\theta,\phi)$ are the spherical
harmonics.\footnote{The factor $\sqrt{2/\pi}$ and the inclusion of $k$ in the expansion coefficient normalization
ensure orthonormality. The complete set of functions $\{j_\ell(kr)Y_{\ell m}(\hat{\mathbf{r}})\}$ forms a basis on
the domain $0\le r<\infty$ and the unit sphere, respectively. This is the natural generalization of plane-wave
Fourier expansion to spherical geometry.}

The inverse transformation, which extracts the SFB coefficients from the field, is:
\begin{equation}
f_{\ell m}(k) = \sqrt{\frac{2}{\pi}}
\int_0^\infty dr\,r^2
\int d\Omega\,
f(\mathbf{r})\,k j_\ell(kr)\,Y^*_{\ell m}(\theta,\phi),
\label{eq:sfb_inverse}
\end{equation}
where $d\Omega = \sin\theta\,d\theta\,d\phi$ is the solid angle element.

\subsubsection{Cartesian Fourier decomposition as reference}

For comparison, the standard Cartesian Fourier transform decomposes the same field as:
\begin{equation}
f(\mathbf{x}) = \frac{1}{(2\pi)^3}
\int d^3\mathbf{k}\,\tilde{f}(\mathbf{k})\,e^{i\mathbf{k}\cdot\mathbf{x}},
\label{eq:cartesian_fourier}
\end{equation}
with inverse:
\begin{equation}
\tilde{f}(\mathbf{k}) = \int d^3\mathbf{x}\,f(\mathbf{x})\,e^{-i\mathbf{k}\cdot\mathbf{x}}.
\end{equation}

The SFB decomposition and Cartesian Fourier are related through the plane-wave expansion in spherical harmonics.
The key advantage of SFB for cosmology is that it naturally respects the spherical symmetry of the observed sky
and the radial structure imposed by redshift-dependent selection functions.\footnote{While Cartesian Fourier is
convenient for periodic boundary conditions (e.g., in simulations), SFB is more natural for surveys where we observe
from a single point (the observer) outward, and where the radial direction has a distinct physical meaning.}

\subsection{The Power Spectrum in Ideal Settings}

\subsubsection{Definition under Statistical Isotropy and Homogeneity}

When a field satisfies the conditions of Statistical Isotropy and Homogeneity (SIH), its power spectrum is defined
through the two-point correlation of the SFB coefficients:
\begin{equation}
\bigl\langle f_{\ell m}(k)\,f^*_{\ell' m'}(k')\bigr\rangle
= C_\ell(k)\,\delta(k-k')\,\delta_{\ell\ell'}\,\delta_{mm'},
\label{eq:sfb_power_def}
\end{equation}
where $\langle\cdot\rangle$ denotes an ensemble average, and $C_\ell(k)$ is the three-dimensional SFB power spectrum.

Remarkably, under the SIH condition, $C_\ell(k)$ is \emph{independent} of the multipole index $\ell$ and
the azimuthal index $m$:\footnote{This independence is a profound consequence of isotropy. Because the field
looks the same in all directions, no particular $\ell$ mode is preferred over any other. The index $\ell$ describes
the angular pattern, but the amplitude of fluctuations depends only on the radial scale $k$, not on how those
fluctuations are oriented in the sky.}
\begin{equation}
C_\ell(k) = P(k) \quad \text{(SIH condition)},
\label{eq:sfb_radial_only}
\end{equation}
where $P(k)$ is simply the three-dimensional isotropic power spectrum, identical to that obtained from Cartesian
Fourier analysis.

Correspondingly, the Cartesian Fourier power spectrum is defined as:
\begin{equation}
\bigl\langle\tilde{f}(\mathbf{k})\,\tilde{f}^*(\mathbf{k}')\bigr\rangle
= (2\pi)^3\,P(k)\,\delta^3(\mathbf{k}-\mathbf{k}'),
\label{eq:cartesian_power_def}
\end{equation}
showing that the two power spectra are identical when the field is isotropic.

\subsubsection{Physical interpretation of radial independence}

The result that $C_\ell(k)$ is independent of $\ell$ under SIH expresses a fundamental fact: in a homogeneous and
isotropic universe, there is no preferred direction. A structure at some wavenumber $k$ is equally likely to be
oriented radially (low $\ell$) or tangentially (high $\ell$) on the sky. This is the basis for compression of the full
3D power spectrum information into a single function $P(k)$.\footnote{This statement applies to the underlying
universe model. In practice, finite survey volumes and the redshift-dependent evolution of structure break SIH,
and $C_\ell(k)$ develops dependence on both $k$ and $\ell$.}

\subsection{Breaking Homogeneity: Finite Survey Selection Functions}

\subsubsection{How selection functions couple the radial and angular dimensions}

Real surveys do not have infinite volume. Instead, they are limited by:
\begin{itemize}
\item A radial selection function $\phi(r)$ that describes the survey footprint in comoving distance (determined
by redshift completeness, photometric redshift uncertainties, etc.).
\item Angular selection, typically accounted for by masking or mode-mixing corrections.
\end{itemize}

When selection is applied, the observed field becomes:
\begin{equation}
f^{\text{obs}}(\mathbf{r}) = \phi(r)\,f(\mathbf{r}),
\label{eq:observed_field_with_selection}
\end{equation}
where $\phi(r)$ is assumed to depend only on radial distance (full-sky coverage is assumed for now).

This multiplication breaks homogeneity in the radial direction.\footnote{Mathematically, the convolution theorem in
Fourier space tells us that multiplying a field by a window function in position space causes mixing of Fourier modes.
This mode mixing is the fundamental reason why observed spectra depend on both $k$ and $\ell$.}

\subsubsection{The observed SFB power spectrum with selection}

For the observed field, the definition of power spectrum generalizes to:
\begin{equation}
\bigl\langle f^{\text{obs}}_{\ell m}(k)\,f^{\text{obs}*}_{\ell' m'}(k')\bigr\rangle
= C^{\text{obs}}_\ell(k,k')\,\delta_{\ell\ell'}\,\delta_{mm'},
\label{eq:observed_power_general}
\end{equation}
where now $C^{\text{obs}}_\ell(k,k')$ is a function of both $k$ and $k'$ (a matrix in Fourier space at each $\ell$,
rather than a single function).

When the underlying field $f(\mathbf{r})$ obeys SIH, the observed spectrum can be written in terms of the true
spectrum $P(k)$ and a window function $W_\ell(k,k')$ that encodes the effect of selection:
\begin{equation}
C^{\text{obs}}_\ell(k,k') = \left(\frac{2}{\pi}\right)^2
\int_0^\infty dk''\, k''^2\,P(k'')\,W_\ell(k,k'')\,W_\ell(k',k''),
\label{eq:observed_power_window}
\end{equation}
where the window function is defined as:
\begin{equation}
W_\ell(k,k') = \int_0^\infty dr\,r^2\,\phi(r)\,j_\ell(kr)\,j_\ell(k'r).
\label{eq:window_function_def}
\end{equation}

This expression is the cornerstone result: \textbf{the observed power spectrum is a convolution of the true power
spectrum with the window function}.\footnote{This is a direct consequence of the convolution theorem applied to
the multiplication $\phi(r)f(\mathbf{r})$ in the SFB basis. The factor $(2/\pi)^2$ is a normalization convention
that depends on the chosen expansion basis.}

\subsection{Special Limits and Geometric Intuition}

\subsubsection{The radial limit: deep survey}

In the limit where the radial selection function is very broad (survey extends over a large redshift range), the
window function becomes sharply peaked on the diagonal $k=k'$. Specifically, if $r_0$ characterizes the survey
depth (roughly the distance to the survey median), and if
\begin{equation}
r_0 k \gg \sqrt{2\ell(\ell+1)},
\label{eq:radialisation_condition}
\end{equation}
then the window function approaches:
\begin{equation}
W_\ell(k,k') \approx \frac{\pi}{2k'}\,\delta(k-k'),
\label{eq:window_radial_limit}
\end{equation}
and thus:
\begin{equation}
C^{\text{obs}}_\ell(k,k') \approx C_\ell(k)\,\delta(k-k') = P(k)\,\delta(k-k').
\label{eq:observed_power_radial_limit}
\end{equation}

This \textbf{radialisation} effect means that for deep surveys, the power spectrum becomes nearly diagonal: each
observed radial mode $k$ couples only to modes with the same $k$ value. The information is compressed into radial
wavenumber, independent of $\ell$.\footnote{Physically, radialisation occurs because in a deep survey, most of the
volume is far from the observer. Structures at different radial distances contribute coherently to the angular pattern
only if they have similar radial wavelengths. This causes mode mixing in $\ell$ to average out.}

\subsubsection{The tangential limit: thin shell}

Conversely, if the selection function is a thin delta function at distance $r_*$:
\begin{equation}
\phi(r) = r_*\,\delta(r-r_*),
\label{eq:thin_shell_selection}
\end{equation}
then the window function becomes:
\begin{equation}
W_\ell(k,k') = r_*^3\,j_\ell(kr_*)\,j_\ell(k'r_*),
\label{eq:window_thin_shell}
\end{equation}
and the observed spectrum is:
\begin{equation}
C^{\text{obs}}_\ell(k,k') = \left(\frac{2}{\pi}\right)^2 kk'\,j_\ell(kr_*)\,j_\ell(k'r_*)
\int_0^\infty dk''\,k''^2\,P(k'')\,\bigl[r_*^3 j_\ell(k''r_*)\bigr]^2.
\label{eq:observed_power_thin_shell}
\end{equation}

This can be related directly to the two-dimensional spherical harmonic power spectrum measured in that shell:
\begin{equation}
C^{\text{2D}}_\ell = \frac{2}{\pi}\int_0^\infty dk\,k^2\,P(k)\,\bigl[W^{\text{2D}}_\ell(k)\bigr]^2,
\label{eq:2d_power_spectrum}
\end{equation}
where:
\begin{equation}
W^{\text{2D}}_\ell(k) = \int_0^\infty dr\,r^2\,\phi(r)\,j_\ell(kr).
\label{eq:2d_window}
\end{equation}

For the thin shell, the observed 3D power becomes:
\begin{equation}
C^{\text{obs}}_\ell(k,k') = \frac{2}{\pi}\,k\,k'\,j_\ell(kr_*)\,j_\ell(k'r_*)\,C^{\text{2D}}_\ell.
\label{eq:connection_2d_3d}
\end{equation}

Thus all the information is compressed into the 2D power spectrum $C^{\text{2D}}_\ell$, which depends only on $\ell$,
not on the individual Fourier modes $k,k'$.\footnote{This makes sense geometrically: if the survey is only one shell,
there is no radial information to extract, only angular. The 3D spectrum collapses to its 2D projection.}

\subsection{Evolution and Bias: Realistic Galaxy Fields}

\subsubsection{The biased and evolving field structure}

In realistic surveys, the galaxy field is not an unbiased tracer of the underlying dark matter distribution. Instead,
the galaxy power spectrum at a given redshift is related to the matter power spectrum by scale-dependent bias and
redshift-dependent growth:
\begin{equation}
P_{\text{gal}}^{\text{evol}}(k,r) = b^2(r)\,D^2(r)\,P_{\text{m}}(k),
\label{eq:evolving_biased_power}
\end{equation}
where:
\begin{itemize}
\item $b(r)$ is the linear galaxy bias factor (assumed scale-independent in the linear regime), which quantifies how
strongly galaxies cluster relative to the underlying matter.
\item $D(r)$ is the linear growth factor, which describes how matter perturbations grow with cosmic time (and thus
with distance along the line of sight, if redshift space is the observable).
\item $P_{\text{m}}(k)$ is the matter power spectrum at the present epoch (a constant reference).
\end{itemize}

\subsubsection{Modified window function with evolution}

When the field has evolution and bias, the window function that relates the observed power to the true matter
power must include these dependencies:
\begin{equation}
W^{\text{evol}}_\ell(k,k') = \int_0^\infty dr\,r^2\,\phi(r)\,D(r)\,b(r)\,j_\ell(kr)\,j_\ell(k'r).
\label{eq:window_evolved}
\end{equation}

The observed galaxy power spectrum then becomes:
\begin{equation}
C^{\text{gal}}_\ell(k,k') = \left(\frac{2}{\pi}\right)^2
\int_0^\infty dk''\,k''^2\,P_{\text{m}}(k'')\,W^{\text{evol}}_\ell(k,k'')\,W^{\text{evol}}_\ell(k',k''),
\label{eq:observed_galaxy_power}
\end{equation}

This shows that evolution and bias effectively get ``baked into'' the window function, so that the convolution integral
structure is preserved.\footnote{The factors $D(r)$ and $b(r)$ vary slowly over the survey volume in many cases (especially
for deep surveys). This means that for some applications, one can approximate them as constants evaluated at a
reference distance, simplifying the analysis.}

\subsubsection{Linear and nonlinear regimes}

The linear bias approximation (Equation \ref{eq:evolving_biased_power}) holds when the wavelength of interest is
much larger than the nonlinear scale $k_{\text{nl}}(z)$. For LSST-like surveys and the scales of interest, this
typically means:
\begin{equation}
k < k_{\text{max}}(z) \approx 0.12\,h\,\text{Mpc}^{-1} \quad \text{at } z=0,
\end{equation}
and:
\begin{equation}
k < k_{\text{max}}(z) \approx 0.25\,h\,\text{Mpc}^{-1} \quad \text{at } z=2.
\end{equation}

Beyond these scales, the nonlinear power spectrum must be used, and the effective bias may become scale-dependent
and more complex to model.\footnote{The nonlinear power spectrum can be computed using $N$-body simulations or
fitted formulae (e.g., Halofit). The scale of mode-coupling and nonlinearity grows with redshift because structures
have had less time to gravitationally collapse.}

\subsection{Window Functions in Practice}

\subsubsection{Gaussian selection function}

A common model for the radial selection function is a Gaussian distribution:
\begin{equation}
\phi(r) = e^{-(r/r_0)^2},
\label{eq:gaussian_selection}
\end{equation}
where $r_0$ is a scale parameter related to the median redshift of the survey.

For this choice, the window function can be evaluated analytically:
\begin{equation}
W_\ell(k,k') = \frac{\pi r_0^2}{4}\sqrt{\frac{k}{k'}}
\exp\left[-r_0^2\frac{k^2+k'^2}{4}\right]
I_{\ell+1/2}\left(\frac{r_0^2 k k'}{2}\right),
\label{eq:window_gaussian}
\end{equation}
where $I_\nu(x)$ is the modified Bessel function of the first kind.\footnote{This analytical form is valuable because
it allows rapid evaluation of the convolution integral. For large arguments, the modified Bessel function can be
computed using asymptotic forms or scaled representations (e.g., $\tilde{I}_\nu(x)=e^{-x}I_\nu(x)$) to avoid numerical
overflow.}

\subsubsection{Radialisation for Gaussian surveys}

The radialisation condition (Equation \ref{eq:radialisation_condition}) can be derived from the Gaussian window
function in the limit $r_0\to\infty$. Using the asymptotic expansion of the modified Bessel function for $x\gg\alpha^2-1/4$:
\begin{equation}
I_\alpha(x) \approx \frac{e^x}{\sqrt{2\pi x}},
\label{eq:bessel_asymptotic}
\end{equation}
one obtains:
\begin{equation}
W_\ell(k,k') \approx \sqrt{\frac{\pi}{4}}\frac{r_0}{k'}
\exp\left[-r_0^2\frac{(k-k')^2}{4}\right]
\quad \text{when } r_0^2 k k' \gg 2\ell(\ell+1).
\label{eq:window_gaussian_radial_limit}
\end{equation}

In this regime, the window becomes a narrowly peaked Gaussian around $k=k'$, and thus:
\begin{equation}
W_\ell(k,k') \approx \frac{\pi}{2k'}\,\delta(k-k'),
\label{eq:window_dirac_limit}
\end{equation}
recovering the result of Equation \ref{eq:window_radial_limit}.\footnote{The condition $r_0^2 k k' \gg 2\ell(\ell+1)$
is equivalent to $r_0 k \gg \sqrt{2\ell(\ell+1)}$, which is the radialisation condition. For $r_0=1400\,h^{-1}\text{Mpc}$
and $\ell=10$, this requires $k \gg 0.05\,h\,\text{Mpc}^{-1}$, while for $\ell=100$, it requires $k\gg 0.16\,h\,\text{Mpc}^{-1}$.
This quantifies the scales at which different angular multipoles become effectively radial.}

\subsection{Application to Baryon Acoustic Oscillations}

\subsubsection{BAO signature in the power spectrum}

The matter power spectrum exhibits oscillatory features due to the imprint of baryon acoustic oscillations (BAOs),
which arise from the interplay of gravity and radiation pressure in the early universe. These oscillations can be
parameterized by a ratio:
\begin{equation}
R^P(k) = \frac{P_{\text{with BAO}}(k)}{P_{\text{smooth}}(k)},
\label{eq:bao_ratio_fourier}
\end{equation}
where $P_{\text{with BAO}}(k)$ includes the wiggles and $P_{\text{smooth}}(k)$ is the power spectrum with oscillations
removed (often computed via spline fitting).

\subsubsection{BAOs in SFB space}

In the SFB framework, one can similarly define a ratio at each $(\ell,k)$ point:
\begin{equation}
R^C_\ell(k) = \frac{C^{\text{obs,with BAO}}_\ell(k,k)}{C^{\text{obs,smooth}}_\ell(k,k)},
\label{eq:bao_ratio_sfb}
\end{equation}
where the diagonal element $C^{\text{obs}}_\ell(k,k)$ is used (cf. Equation \ref{eq:observed_power_general}).

For shallow surveys, the BAO wiggles depend on both $k$ and $\ell$: low-$\ell$ modes contain information about
the BAO scale measured tangentially (through the angular direction on the sky), while high-$k$ modes contain
radial BAO information.

For deep surveys, the radialisation phenomenon causes the BAO signal to become increasingly concentrated in the
radial modes (the diagonal $C^{\text{obs}}_\ell(k,k)$), with weaker angular dependence. This is because the large
survey volume allows coherent mode cancellation of high-$\ell$ modes, leaving only the radial BAO signature.\footnote{This
is a key physical insight: deep surveys are naturally sensitive to radial BAO features, while shallow surveys probe
both radial and angular BAO signatures. Combining information across different $\ell$ at a fixed $k$ can help reduce
cosmic variance and improve constraints on the BAO scale.}

\subsubsection{Convergence to Fourier space}

In the limit of infinite radial coverage (or equivalently, $r_0\to\infty$), the SFB power spectrum converges to
the Cartesian Fourier power spectrum:
\begin{equation}
\lim_{r_0\to\infty} C^{\text{obs}}_\ell(k,k) = P(k),
\label{eq:convergence_fourier}
\end{equation}
and thus:
\begin{equation}
\lim_{r_0\to\infty} R^C_\ell(k) = R^P(k).
\label{eq:convergence_bao_ratio}
\end{equation}

This limit justifies the use of the standard Fourier-space BAO analysis for very deep surveys, and provides a check
on the SFB formalism.\footnote{The rate of convergence depends on the angular scale: high-$\ell$ modes converge
faster to the Fourier limit than low-$\ell$ modes, because the radialisation condition is satisfied for larger $\ell$
at fixed $k$.}

\subsection{Summary of Key Results}

\begin{itemize}

\item \textbf{SFB decomposition} provides a natural basis for 3D fields on the sky, with complete orthonormality
and efficient handling of selection functions.

\item \textbf{Isotropic power spectrum}: Under SIH conditions, $C_\ell(k)=P(k)$ is independent of $\ell$, recovering
the standard 3D isotropic spectrum.

\item \textbf{Selection-induced mode mixing}: Finite survey volumes break SIH and cause the observed spectrum
$C^{\text{obs}}_\ell(k,k')$ to depend on both radial modes $k$ and $k'$. This mixing is encoded in the window function
$W_\ell(k,k')$.

\item \textbf{Radialisation}: For deep surveys with large $r_0$, the condition $r_0 k \gg \sqrt{2\ell(\ell+1)}$ is
satisfied, and the spectrum becomes diagonal. High-$\ell$ modes radialise at smaller $k$ than low-$\ell$ modes.

\item \textbf{Geometric limits}: The thin-shell limit recovers 2D spherical harmonic spectra, while the deep-survey
limit recovers 1D Fourier spectra, demonstrating that SFB unifies these limiting cases.

\item \textbf{BAO applications}: BAO wiggles manifest in both $k$ and $\ell$ dimensions in the SFB spectrum. Radialisation
causes wiggles to concentrate in the radial direction for deep surveys, consistent with the Fourier BAO picture.

\end{itemize}
